\documentclass[11pt]{article}

\usepackage[margin=1in]{geometry}
\usepackage{amsmath, amssymb, amsthm}
\usepackage{array}
\usepackage{booktabs}
\usepackage{enumitem}
\usepackage{hyperref}

\newcolumntype{L}[1]{>{\raggedright\arraybackslash}p{#1}}

\title{First Two-Class Appointment Simulation: Recommended Formulation}
\author{Alexandre SEPULVEDA de DIETRICH\\Under the supervision of Pr.~Carri Chan}
\date{\today}

\begin{document}
\maketitle

\section{Purpose}

This note defines the first discrete-event appointment simulation used in this project and
clarifies the notation so it matches the underlying dynamics of the problem.
The main design goal is to keep the first simulator simple enough to run, inspect, and extend,
while using a formulation that remains valid once we add richer allocation policies,
heterogeneous patient behavior, or optimization layers.

The note aligns the simulator with the queueing models in:
\begin{itemize}[nosep]
    \item Green and Savin (2008), \emph{Reducing Delays for Medical Appointments: A Queueing Approach},
    \item Liu and Ziya (2014), \emph{Panel Size and Overbooking Decisions for Appointment-based Services under Patient No-shows},
    \item Liu (2015), \emph{Optimal Choice for Appointment Scheduling Window under Patient No-Show Behavior}.
\end{itemize}

All three papers study the same structural mechanism:
\begin{quote}
more backlog $\Rightarrow$ more delay $\Rightarrow$ worse attendance $\Rightarrow$ lower effective capacity.
\end{quote}

The simulation keeps that mechanism, but implements it on an explicit rolling calendar with
two patient classes, explicit cancellations, and a pluggable slot-allocation rule.

\section{Why the Naive Formulation is Not Sufficient}

An initial formulation might write
\[
X_{i,\tau}^{D} = \text{number of class-$i$ patients booked for day } D+\tau
\]
and then also use $\tau$ inside $b_i(\tau)$ and $\xi_i(\tau)$.
That notation is not the best formulation for this problem.

\subsection{Problem 1: It Overloads the Delay Variable}

There are two different delays in the system:
\begin{enumerate}[nosep]
    \item the \emph{remaining delay} from the current day to the appointment,
    \item the \emph{original delay seen at booking time}.
\end{enumerate}

These play different roles.
The remaining delay is a state variable.
The original delay at booking is the natural argument of the behavioral functions.

If the same symbol $\tau$ is used for both, then the model becomes ambiguous.
In particular, it can incorrectly suggest that no-show probability should be recomputed from
remaining delay as the appointment approaches.
That is not the intended behavioral mechanism here.

\subsection{Problem 2: The Aggregate Day-Level Count is Not the True State}

The vector $\{X_{i,\tau}^{D}\}$ is a useful summary, but it is not the full simulator state once
we include:
\begin{itemize}[nosep]
    \item slot-level capacity,
    \item day-by-day cancellations that reopen specific future slots,
    \item alternative booking policies that may choose different open slots within the same day.
\end{itemize}

For those features, the actual state must retain the slot-level calendar.

\subsection{Problem 3: It Hides the Link to the Literature}

In the source papers, attendance depends on the backlog or delay the patient sees when the
appointment is accepted.
The relevant quantity is therefore the \emph{booking-time delay}, not the residual delay just
before service.

The simulator should preserve that same logic in explicit-calendar form.

\section{Recommended Formulation}

\subsection{Decision Epochs}

Let time advance in slot-sized decision epochs.
Write
\[
t = (D,s),
\]
where:
\begin{itemize}[nosep]
    \item $D$ is the calendar day,
    \item $s \in \{0,1,\dots,S-1\}$ is the active slot within day $D$.
\end{itemize}

In the first version of the simulator,
\[
H = 15 \text{ days}, \qquad S = 25 \text{ slots per day},
\]
where $H$ is the booking horizon.

\subsection{True State: Slot-Level Calendar}

The true simulator state at time $t=(D,s)$ is the rolling slot-level calendar.
For each residual day offset
\[
r \in \{0,1,\dots,H-1\}
\]
and each slot index
\[
m \in \{0,1,\dots,S-1\},
\]
define
\[
Y_t(r,m).
\]

Interpretation:
\begin{itemize}[nosep]
    \item $Y_t(r,m)=0$ if slot $(D+r,m)$ is currently open,
    \item $Y_t(r,m)=\left(i,\tau^{\mathrm{book}},c\right)$ if slot $(D+r,m)$ is occupied by a class-$i$ patient whose original booking delay was $\tau^{\mathrm{book}}$, with cancellation marker $c$.
\end{itemize}

The cancellation marker $c$ is either:
\begin{itemize}[nosep]
    \item a future calendar day on which the appointment will be canceled before service, or
    \item $\varnothing$ if the appointment will not pre-cancel.
\end{itemize}

This is the minimal state representation that supports:
\begin{itemize}[nosep]
    \item explicit holes created by cancellation,
    \item slot-level FCFS or non-FCFS allocation,
    \item class-specific booked-delay tracking,
    \item later policy experiments that depend on slot position.
\end{itemize}

\subsection{Derived Day-Start Summary}

Although the slot-level calendar is the true state, it is still useful to retain an aggregate
day-start summary that resembles the literature.
At the start of day $D$, after applying cancellations for day $D$ and before processing slot $0$,
define
\[
X_{i,r}^{D}
=
\sum_{m=0}^{S-1}
\mathbf{1}\{Y_{(D,0)}(r,m)\text{ is occupied by a class-}i\text{ patient}\},
\]
for
\[
i \in \{1,2\}, \qquad r \in \{0,1,\dots,H-1\}.
\]

Interpretation:
\[
X_{i,r}^{D}
=
\text{number of class-}i\text{ patients scheduled at the start of day }D\text{ for day }D+r.
\]

This is the correct interpretation of the state-count variable.
The index $r$ is a \emph{residual delay}, not a booking-time delay.

\subsection{Booking-Time Delay}

Now define the behavioral delay variable separately:
\[
\tau^{\mathrm{book}}
=
\text{delay in days between the booking day and the appointment day offered to the patient.}
\]

If a patient arrives on day $d$ and is offered a slot on day $d+\tau^{\mathrm{book}}$, then:
\begin{itemize}[nosep]
    \item $b_i(\tau^{\mathrm{book}})$ is the probability that class-$i$ declines the offer,
    \item if the patient accepts, $\tau^{\mathrm{book}}$ is stored with the booking,
    \item $\xi_i(\tau^{\mathrm{book}})$ is the no-show probability used at appointment time if the patient remains scheduled.
\end{itemize}

This separation is the key modeling choice.
Residual delay belongs to the state.
Booking delay belongs to the behavioral rules.

\section{Patient Classes}

Each class $i \in \{1,2\}$ is described by
\[
C_i = \{\lambda_i, b_i(\tau^{\mathrm{book}}), \phi_i, \xi_i(\tau^{\mathrm{book}})\}.
\]

The interpretation is:
\begin{itemize}[nosep]
    \item $\lambda_i$: expected arrival rate of class-$i$ requests per slot,
    \item $b_i(\tau^{\mathrm{book}})$: probability the patient declines the offered slot because the offered delay is too large,
    \item $\phi_i$: one-shot pre-appointment cancellation probability after booking,
    \item $\xi_i(\tau^{\mathrm{book}})$: no-show probability at appointment time, conditional on the patient still being scheduled.
\end{itemize}

In the current first version, $\phi_i$ is constant within class.
Later versions may generalize this to $\phi_i(\tau^{\mathrm{book}})$ if cancellation is also delay-sensitive.

\section{Event Dynamics}

\subsection{Slot-Level Event Order}

At each active slot $(D,s)$, the simulator executes the following steps.

\begin{enumerate}[leftmargin=*, itemsep=0.35em]
    \item Apply all cancellations whose cancellation day is $D$.
    \item Generate arrivals
    \[
    A_i(D,s) \sim \mathrm{Poisson}(\lambda_i), \qquad i \in \{1,2\}.
    \]
    \item Randomize the order of individual arrivals within the slot.
    \item For each arrival, call the allocation policy to select an open future slot.
    \item If a slot on day $D+r$ is offered, interpret
    \[
    \tau^{\mathrm{book}} = r.
    \]
    \item The patient accepts with probability
    \[
    1 - b_i(\tau^{\mathrm{book}}).
    \]
    \item If the patient books and $\tau^{\mathrm{book}} \ge 1$, then with probability $\phi_i$ a cancellation day is sampled uniformly from the days strictly before the appointment day.
    \item When the active slot is served, the patient in that slot, if any, attends with probability
    \[
    1 - \xi_i(\tau^{\mathrm{book}}),
    \]
    where $\tau^{\mathrm{book}}$ is the original delay stored when the patient booked.
\end{enumerate}

\subsection{FCFS Baseline Policy}

In the baseline version, the allocation rule is:
\begin{quote}
choose the earliest open slot strictly after the active slot.
\end{quote}

That rule is first-come first-served in calendar order, but the simulator is intentionally
written so the slot-selection rule can later be replaced without changing the event loop.
The current codebase now includes two additional initial-model examples:
\begin{itemize}[nosep]
    \item a strict reserved-capacity policy that protects a specified fraction or count of daily slots for each class,
    \item a class-specific booking-window policy that restricts class $i$ to delays no larger than $k_i$.
\end{itemize}

\section{Performance Measures}

\subsection{Arrival-Cohort Measures}

For patients who arrive during the measured window, define:
\begin{align*}
\mathrm{Arrivals}_i &= \text{number of class-}i\text{ appointment requests}, \\
\mathrm{Booked}_i &= \text{number of class-}i\text{ requests that accept an offer}, \\
\mathrm{Balked}_i &= \text{number of class-}i\text{ requests that do not book}, \\
\mathrm{Canceled}_i &= \text{number of class-}i\text{ booked patients who pre-cancel}, \\
\mathrm{NoShow}_i &= \text{number of class-}i\text{ booked patients who remain scheduled but do not attend}, \\
\mathrm{Served}_i &= \text{number of class-}i\text{ booked patients who attend}, \\
\bar{\tau}^{\mathrm{book}}_i &= \text{mean original booking delay among class-}i\text{ booked patients}.
\end{align*}

The main service metric in the first version is
\[
\mathrm{ServiceRate}_i
=
\frac{\mathrm{Served}_i}{\mathrm{Booked}_i}.
\]

Aggregate values are formed by summing over classes and taking the corresponding weighted mean.

\subsection{Slot-Based Measures}

Over measured service slots, define:
\begin{align*}
\mathrm{BookedSlots} &= \text{slots with a booked patient at service time}, \\
\mathrm{ServedSlots} &= \text{slots where the booked patient attends}, \\
\mathrm{NoShowSlots} &= \text{slots where a booked patient is present on the schedule but does not attend}, \\
\mathrm{EmptySlots} &= \text{slots with nobody scheduled at service time}.
\end{align*}

Two utilization measures are reported:
\[
\rho^{\mathrm{booked}}
=
\frac{\mathrm{BookedSlots}}{HS_{\mathrm{measured}}},
\qquad
\rho^{\mathrm{attended}}
=
\frac{\mathrm{ServedSlots}}{HS_{\mathrm{measured}}}.
\]

The first measures scheduled load.
The second measures realized attended load.

\section{Comparison to the Literature}

\subsection{Common Structural Mapping}

The source papers use more stylized queue states than the simulator, but the economic and
behavioral interpretation is the same.
In each paper, attendance depends on the congestion level a patient sees when entering the
system.

In the analytical papers, that congestion variable is usually:
\begin{itemize}[nosep]
    \item queue length,
    \item backlog level,
    \item or the number of appointments already ahead in the system.
\end{itemize}

In the simulator, the explicit-calendar analog is the booking-time delay $\tau^{\mathrm{book}}$.
That is why $\tau^{\mathrm{book}}$, rather than residual delay $r$, should drive
$b_i(\cdot)$ and $\xi_i(\cdot)$.

\subsection{Green and Savin (2008)}

Green and Savin study a single-server appointment queue in which no-show probability increases
with the backlog seen when the appointment is made.
Their no-show curve is of the exponential-saturation form
\[
\gamma(k) = \gamma_{\max} - (\gamma_{\max} - \gamma_0)e^{-k/C},
\]
where $k$ is a backlog measure.

The mapping to the simulator is:
\begin{itemize}[nosep]
    \item their backlog state corresponds to our explicit slot calendar,
    \item their delay-sensitive no-show mechanism corresponds to our $\xi_i(\tau^{\mathrm{book}})$,
    \item their stylized queue-length state is refined into the explicit schedule $Y_t(r,m)$.
\end{itemize}

The simulation differs from Green and Savin in several ways:
\begin{itemize}[nosep]
    \item it allows two patient classes rather than one homogeneous stream,
    \item it includes explicit balking before booking,
    \item it includes explicit pre-appointment cancellation timing,
    \item it does not yet model no-show rescheduling.
\end{itemize}

The most important conceptual similarity is this:
attendance deteriorates with the delay the patient faced when entering the schedule.

\subsection{Liu and Ziya (2014)}

Liu and Ziya study panel size and overbooking decisions in stylized queueing models with
delay-dependent attendance.
Their state variable is again a queue-length-type object, and a patient's show-up probability
depends on the congestion level observed at booking time.

Relative to Liu and Ziya, the simulator:
\begin{itemize}[nosep]
    \item fixes daily slot capacity at $25$ in the first version,
    \item does not yet optimize overbooking or panel size,
    \item keeps the booking calendar explicit rather than compressing it into a single queue-length state,
    \item is designed to allow class-specific balking, cancellation, and attendance curves.
\end{itemize}

Their model is therefore more compact and optimization-oriented.
Our simulator is more descriptive and state-rich.
The key shared mechanism remains the same:
\[
\text{booking-time congestion} \longrightarrow \text{attendance}.
\]

\subsection{Liu (2015)}

Liu studies the optimal appointment-window size $K$ under delay-dependent no-shows.
That paper's main control is the admissible backlog or booking window.

The simulator's rolling horizon length $H$ plays the same operational role as a hard scheduling
window, but here the horizon is fixed in the baseline implementation.
This makes the simulator a natural platform for later experiments in which:
\begin{itemize}[nosep]
    \item $H$ is varied,
    \item different allocation policies are compared,
    \item or class-specific admission rules are introduced.
\end{itemize}

Again, the important mapping is not the exact state variable but the interpretation:
the patient's future attendance depends on the delay accepted at booking time.

\subsection{Side-by-Side Comparison}

\begin{center}
\small
\begin{tabular}{L{0.18\textwidth} L{0.20\textwidth} L{0.22\textwidth} L{0.24\textwidth}}
\toprule
Model / paper & State abstraction & Delay variable that drives attendance & Main difference from this simulator \\
\midrule
Green and Savin (2008) & queue length / backlog in an appointment queue & backlog seen when appointment is made & analytical single-class queue; no explicit slot calendar \\
\addlinespace
Liu and Ziya (2014) & queue length in $M/M/1$ or $M/D/1$ style models & congestion at booking time & optimization over panel size and overbooking, not explicit calendar simulation \\
\addlinespace
Liu (2015) & queue/window state with finite scheduling window $K$ & delay implied by admissible backlog at booking & optimization over window size rather than fixed-horizon simulation \\
\addlinespace
This simulator & explicit rolling calendar $Y_t(r,m)$ plus derived $X_{i,r}^D$ & original booking delay $\tau^{\mathrm{book}}$ & richer operational detail, two classes, explicit cancellation holes, policy plug points \\
\bottomrule
\end{tabular}
\normalsize
\end{center}

\section{Why this Formulation is the Most Appropriate for This Problem}

This formulation is the most appropriate starting point for the present project for four reasons.

\subsection{1. It Matches the Behavioral Mechanism}

The patient decides whether to book when seeing an offered delay.
If they book, their future no-show behavior should be tied to that accepted delay.
That is exactly what $\tau^{\mathrm{book}}$ captures.

\subsection{2. It Preserves the Literature's Logic While Using an Explicit Calendar}

The source papers compress the calendar into queue length for tractability.
The simulator does not need to do that.
It can retain the slot-level calendar and still summarize the system in a literature-friendly way
through $X_{i,r}^{D}$.

\subsection{3. It Is Markov for the Planned Extensions}

Once cancellations create holes and policies can choose among open slots, a pure day-level count
is no longer enough.
The calendar state $Y_t(r,m)$ is the right object to carry forward if the project later adds:
\begin{itemize}[nosep]
    \item priority rules,
    \item reserved capacity by class,
    \item different appointment windows by class,
    \item overbooking,
    \item dynamic acceptance or rejection decisions.
\end{itemize}

\subsection{4. It Cleanly Separates State and Behavior}

The formulation distinguishes:
\begin{itemize}[nosep]
    \item \emph{where the system is now}: represented by residual delay $r$ in the state,
    \item \emph{what delay the patient accepted}: represented by $\tau^{\mathrm{book}}$ in the behavior functions.
\end{itemize}

That separation is the cleanest way to avoid conceptual mistakes in both code and analysis.

\section{Implementation Mapping}

The current codebase follows this notation directly:
\begin{itemize}[nosep]
    \item the slot calendar stores each patient's class, appointment location, booking day, and $\tau^{\mathrm{book}}$,
    \item the arrival log records the offered delay as \texttt{tau\_offered},
    \item the cohort and slot logs record the stored booking delay as \texttt{tau\_booked},
    \item the day-start state log records the residual day offset as \texttt{residual\_delay}.
\end{itemize}

This is exactly the split recommended above.
In addition, the simulator now exposes a daily journal with access, booking, cancellation,
attendance, and utilization measures, and reports the access metric
\[
\mathbb{P}\!\left(\tau^{\mathrm{book}} < T_{\mathrm{access}}\right)
\]
for a configurable target such as $T_{\mathrm{access}} = 30$ days.

\section{References}

\begin{itemize}[leftmargin=*, itemsep=0.25em]
    \item Green, L. V., and S. Savin. 2008. Reducing delays for medical appointments: A queueing approach. \emph{Operations Research} 56(6):1526--1538.
    \item Liu, N. 2015. Optimal choice for appointment scheduling window under patient no-show behavior. \emph{Production and Operations Management} 25(1):128--142.
    \item Liu, N., and S. Ziya. 2014. Panel size and overbooking decisions for appointment-based services under patient no-shows. \emph{Production and Operations Management} 23(12):2209--2223.
\end{itemize}

\end{document}
